\documentclass[10pt,a4paper]{amsart}
\usepackage[margin=19mm]{geometry}
\usepackage{amsmath,amssymb,mathtools}
\usepackage[scr=rsfs,cal=euler]{mathalfa}
\usepackage{xfrac}
\usepackage[unicode,hidelinks,bookmarks=false]{hyperref}
\newcommand{\ie}{i.e.\ }
\newcommand{\cf}{cf.\ }
\newcommand{\resp}{respectively\ }
\newcommand{\Subsection}{\subsection}
\providecommand{\nobreakdash}{\nobreak\hbox{-}}
\setlength{\emergencystretch}{2em}
\multlinegap0pt
\usepackage{amssymb}
\usepackage{tikz}
\newtheorem{proposition}{Proposition}
\newcommand{\C}{\mathbb C}
\newcommand{\Hp}{{\mathbb H}^2}
\newcommand{\R}{\mathbb R}
\title{Bending, entropy and proper affine actions of surface groups}
\author{Martin Bridgeman, Richard Canary, Andres Sambarino}
\date{Source: arXiv:2602.20146v2 (CC BY 4.0)}
\begin{document}
\maketitle

\noindent\emph{Source and licence.} Bending, entropy and proper affine actions of surface groups. Version arXiv:2602.20146v2 (CC BY 4.0), distributed by arXiv under the Creative Commons Attribution 4.0 International licence, http://creativecommons.org/licenses/by/4.0/, as recorded in the arXiv licence metadata for this exact version and shown on the abstract page https://arxiv.org/abs/2602.20146v2. Full licence notice: \texttt{licenses/arxiv-2602.20146-CC-BY-4.0.txt}.\par\smallskip

\noindent\textit{Editorial scope.} Counted unit: the article's proposition theta-mod (source0.tex lines 1695--1697). The article's complete original proof of it is delivered with it (source0.tex lines 1699--1733, one \texttt{proof} environment of 35 lines). No mathematics is written, repaired or re-derived: the statement and the proof are the article's own lines, in the article's own notation, and no retained line is substituted or re-flowed.  No further source-local context is needed: every cross-reference of every retained line resolves inside the two delivered spans.  One declared adaptation: exactly 1 ancillary strings inside the retained spans are omitted (image-inclusion commands and, where the source repeats a label, the repeated label definitions) and no other line differs from the source; the figure environments, with their placement, captions and labels, are kept, and the package records the omitted strings in fidelity receipts bound to the affected bodies.  No retained line contains a citation, so the wrapper carries no bibliography and no bibliography entry.  The retained lines contain the article's own diagram code, which the wrapper typesets with the standard tikz packages; no image file is delivered and the package declares no asset dependency.  Editorial formatting only, in addition: an AMS \texttt{amsart} preamble that restates the article's own declarations and macros used by the retained lines, namely the environments \texttt{proposition} on separate counters, together with the standard packages those lines need.  The retained environments print in this document as Proposition 1 in order of appearance, whereas the article prints the same items with the numbering of its own full document.  The complete native source of the article as distributed in the arXiv e-print for this exact version is bundled unchanged as \texttt{sources/195163/source0.tex}.
\begin{proposition}\label{theta-mod}
If $\mu\in \mathcal{ML}(\Hp)$ is $\theta$-bounded for some $\theta <\frac{\pi}{2}$, then $\Omega_\mu$ is moderately bent.
\end{proposition}

\begin{proof}
Let $(x,y)$ be a bending pair for $\Omega_\mu$. We may assume that  $(x,y) =(0,\infty)$ and that the plane $Q$ lying above the real axis $\R$
is a support plane for the pleated plane $P_\mu$. We can assume that $\Omega_\mu$ contains the upper half-plane.  
Let $P_-$  and $P_+$ be the left and right half-planes for the geodesic $\overline{0\infty}$.
(Notice that if $g$ is not an atom for the measure, then $Q=P_-=P_+$.)

Suppose that $u\in \partial \Omega_\mu\setminus\{0,\infty\}$ lies to the right of $0$.
Let $\eta:\R\to \mathbb H^3$ be the unique geodesic ray so that $\eta(0)=p=(0,0,|u|)\in\overline{0\infty}$ and
$\lim_{t\to +\infty} \eta(t)=u$. Notice that $\eta(\R)$ is perpendicular to $\overline{0\infty}$. 
Let $\gamma:\R\to P_\mu$ be a geodesic  in the intrinsic metric on $P_\mu$ so that  $p=\gamma(0)$  and  $\lim_{t\to +\infty} \gamma(t)=u$. 
Let $\alpha_t:\R\to \mathbb H^3$ be a geodesic  in $\mathbb H^3$ so that  $p=\alpha_t(0)$ and $\alpha_t(d(p,\gamma(t)))=\gamma(t)$.
Since $\mu$ is $\theta$-bounded, 
the angle  $\theta_\mu(\gamma(t),\gamma(0))$ is at most $\theta$. By definition, $\theta_\mu(\gamma(t),\gamma(0))$ is the angle  between 
$-\alpha_t'(0)$  and $-\gamma_-'(0)$, or equivalently,  the angle  between $\alpha_t'(0)$  and $\gamma_-'(0)$.

Since $f_\mu$ extends continuously to an embedding 
$\bar f_\mu:\partial\mathbb H^2\to\partial \mathbb H^3$,  we see that $\alpha_t$ converges to $\eta$, so 
$\alpha_t'(0)$ converges to $\eta'(0)$. Therefore, the angle between $\eta'(0)$ and  $\gamma_-'(0)$  is at most $\theta$. 

Notice that $\gamma_-'(0)$ lies in the portion of $P_-$ which is to the right of $\overline{0\infty}$. We see, by projecting $\eta$ onto the plane that the
Euclidean line segment $R_u$ joining 0 to u makes an angle of at most $\theta$ with portion $R_-$ of $\partial P_-$ which lies to the right of $0$.
Notice that $R_-$ must be a line segment with positive slope, since otherwise the  bending at $\overline{0\infty}$ would be at least $\frac{\pi}{2}$, which would contradict  the fact that
$\mu$ is $\theta$-bounded. Since $R_u$ lies in the lower half plane and makes an angle less than $\frac{\pi}{2}$ with $R_-$ it must have negative slope.
Therefore, $\Re(u)>0$ (see Figure \ref{fig:theta}).

We may similarly show that if $u \in \partial\Omega_\mu$ lies to the left of $0$, then $\Re(u)<0$.
Therefore, if $L$ is the imaginary axis, then $L$ is a round circle in $\hat\C$ which intersects $\partial\Omega_\mu$ transversely at the points
$0$ and $\infty$ and intersects no other point of $\partial\Omega_\mu$.
\begin{figure}[htbp] %  figure placement: here, top, bottom, or page
   \centering
    
   \caption{Plane $L$}
   \label{fig:theta}
\end{figure}
\end{proof}

\end{document}
